% Generated. Do not edit manually. Edit resume_content.json instead.
\renewcommand{\ResumeSummary}{
Ph.D. mechanical engineer combining HVAC product development, CFD/FEA, scientific programming, and prototype testing. At Rheem, developed and experimentally evaluated fan, air-handling, and condensate-management designs, including a fan-shroud prototype demonstrating more than 80\% higher airflow efficiency (CFM/W) than the previous design. Additional experience includes real-time C/C++ HVAC plant modeling and coupled thermal-fluid simulation.
}

\renewcommand{\ResumeSkills}{
\textbf{Thermal and Fluids:} HVAC airflow systems, Heat transfer, CFD, Transient thermal modeling, Convective thermal-fluid coupling\par
\textbf{Mechanical Design and Analysis:} FEA, Structural optimization, Parametric CAD, Design for manufacturability, Assembly and clearance analysis\par
\textbf{Prototypes and Validation:} 3D printing, Airflow testing, Variable-geometry test fixtures, Experimental modal testing, Model calibration\par
\textbf{Programming and CAD:} Python, C/C++, CadQuery, STEP geometry generation, Custom numerical solvers, Engineering workflow automation\par
}

\renewcommand{\RheemBullets}{
\item \textbf{Outdoor Fan-Shroud Optimization}
\begin{itemize}[leftmargin=1.3em,topsep=1pt,itemsep=1pt]
\item Used CFD to optimize shroud geometry and improve airflow distribution across the condenser coil. Fabricated and tested a prototype, demonstrating more than 80\% improvement in airflow efficiency (CFM/W) relative to the previous design.
\end{itemize}
\item \textbf{Airflow-Induced Condensate Management Invention}
\begin{itemize}[leftmargin=1.3em,topsep=1pt,itemsep=1pt]
\item Invented a passive condensate entrainment and slinging system that uses fan-generated pressure differences to draw condensed water into the fan airflow path and improve condensate slinging performance.
\item Performed CFD simulations to investigate airflow behavior, pressure distribution, and condensate entrainment mechanisms.
\item Designed, fabricated, and iteratively refined 3D-printed prototypes to demonstrate the mechanism.
\item Conducted experimental testing to validate condensate entrainment, slinging effectiveness, and overall system performance.
\end{itemize}
\item \textbf{VPAK Air-Handling Unit Development}
\begin{itemize}[leftmargin=1.3em,topsep=1pt,itemsep=1pt]
\item Designed an air-handling unit that delivered approximately 25\% higher airflow than the existing configuration. Fabricated and tested 3D-printed prototypes to assess airflow and refine design assumptions.
\end{itemize}
\item \textbf{Reconfigurable Fan Test Fixture}
\begin{itemize}[leftmargin=1.3em,topsep=1pt,itemsep=1pt]
\item Designed and fabricated a modular, 3D-printed cross-flow fan test fixture with adjustable tongue and scroll geometry and clearances, enabling multiple configurations to be tested without rebuilding the full housing.
\end{itemize}
\item \textbf{Multi-Material Product FEA Analysis}
\begin{itemize}[leftmargin=1.3em,topsep=1pt,itemsep=1pt]
\item Performed finite element analysis (FEA) to support product designs using ABS and ASA plastics, aluminum and steel sheet metal, and expanded polypropylene (EPP) at different densities.
\end{itemize}
\item \textbf{Structural FEA and Assembly Optimization of Modular AHU}
\begin{itemize}[leftmargin=1.3em,topsep=1pt,itemsep=1pt]
\item Performed comprehensive finite element analysis (FEA) of the nine-segment modular AHU assembly to evaluate structural integrity, press-fit connections, and load-bearing capacity. Analyzed deformation under fan and motor loads to maintain critical clearances between the rotating fan lip and stationary bellmouth, preventing mechanical interference and supporting reliable prototype operation.
\end{itemize}
\item \textbf{Invention Disclosures}
\begin{itemize}[leftmargin=1.3em,topsep=1pt,itemsep=1pt]
\item Submitted four patent disclosures covering HVAC innovations, including fan shroud optimization, condensate management, and experimental testing systems.
\end{itemize}
\item \textbf{Modular VPAK Air-Handling Unit Prototype}
\begin{itemize}[leftmargin=1.3em,topsep=1pt,itemsep=1pt]
\item Designed and developed a nine-segment modular VPAK air-handling unit (AHU) prototype, allowing individual sections to be modified independently for airflow optimization and rapid accommodation of design changes, including heater, coil, and component configurations. Used low-cost 3D printing to accelerate design iteration while maintaining a unified, single-piece AHU configuration for production.
\end{itemize}
\item \textbf{Parametric Fan-Shroud Generation}
\begin{itemize}[leftmargin=1.3em,topsep=1pt,itemsep=1pt]
\item Developed a Python/CadQuery-based tool that converts fan diameter, coil geometry, component locations, and design clearances into 3D shroud geometry and STEP files, reducing repetitive CAD modeling and accelerating design iteration across product configurations.
\end{itemize}
\item \textbf{Motor-Mount Structural Optimization}
\begin{itemize}[leftmargin=1.3em,topsep=1pt,itemsep=1pt]
\item Redesigned a motor-mount structure using FEA, reducing material usage by more than 50\% while maintaining stiffness and vibration requirements. Simplified injection-molded components to reduce secondary manufacturing operations and assembly complexity.
\end{itemize}
\item \textbf{Cross-Flow Fan Aerodynamics}
\begin{itemize}[leftmargin=1.3em,topsep=1pt,itemsep=1pt]
\item Used CFD to investigate tongue-vortex interaction, recirculation, eccentric vortex behavior, and the sensitivity of fan performance to scroll and housing geometry.
\end{itemize}
\item \textbf{Fan-Induced Vibration and Modal Analysis}
\begin{itemize}[leftmargin=1.3em,topsep=1pt,itemsep=1pt]
\item Performed FEA-based modal analysis to identify natural frequencies of product components and assess resonance risk by comparing them with fan rotational frequency (RPM/60) and blade-passing frequency.
\end{itemize}
}

\renewcommand{\LTTSBullets}{
\item \textbf{Real-Time C/C++ HVAC Plant Model for SiL and HiL}
\begin{itemize}[leftmargin=1.3em,topsep=1pt,itemsep=1pt]
\item Developed a real-time C/C++ executable HVAC plant model for Software-in-the-Loop (SiL) and Hardware-in-the-Loop (HiL) applications, combining finite element solution of transient heat transfer and velocity-pressure-coupled Navier-Stokes equations with lumped-capacitance heater/cooler models, convective thermal-fluid coupling, and a proportional-derivative (PD) HVAC controller model.
\end{itemize}
\item \textbf{Modular GUI for Aerodynamic Data Analysis}
\begin{itemize}[leftmargin=1.3em,topsep=1pt,itemsep=1pt]
\item Developed a modular GUI for aerodynamic data analysis with well-defined input/output interfaces, improving maintainability and reducing computational and memory overhead.
\end{itemize}
}

\renewcommand{\UTSABullets}{
\item \textbf{Custom 3D Multiphysics Finite Element Solvers}
\begin{itemize}[leftmargin=1.3em,topsep=1pt,itemsep=1pt]
\item Developed custom three-dimensional finite element solvers from first principles for structural mechanics, transient heat transfer, and fluid dynamics, including coupled structural-thermal habitat models and interior thermal-fluid simulations.
\end{itemize}
\item \textbf{NASA Modular-Coupled Virtual Testbed (MCVT)}
\begin{itemize}[leftmargin=1.3em,topsep=1pt,itemsep=1pt]
\item Developed and integrated structural, thermal, and fluid plant models for NASA's Modular-Coupled Virtual Testbed (MCVT), a system-of-systems simulation environment for virtual testing of resilient extraterrestrial habitats.
\end{itemize}
\item \textbf{Thermal Protective Layer Topology Optimization}
\begin{itemize}[leftmargin=1.3em,topsep=1pt,itemsep=1pt]
\item Investigated thermo-mechanical topology optimization of lunar habitat thermal protective layers under different environmental conditions and lunar latitudes.
\end{itemize}
\item \textbf{Experimental Structural Digital Twin and Model Updating}
\begin{itemize}[leftmargin=1.3em,topsep=1pt,itemsep=1pt]
\item Created a structural digital twin of an experimental habitat by acquiring impact-hammer and accelerometer data, identifying natural frequencies and mode shapes, using Crypton optical geometry measurements, and calibrating a parameterized beam-element model to match identified modal properties and measured acceleration histories.
\end{itemize}
}

\renewcommand{\OtherExperience}{

}

\renewcommand{\Demonstrations}{
HVAC duct CFD: \url{https://www.linkedin.com/feed/update/urn:li:activity:7258219618144538624/}\par
Lunar habitat thermal topology optimization: \url{https://youtu.be/EGVZgg1pb6Q}\par
MCVT structural-mechanical model: \url{https://youtu.be/TAA66P_UXIM}\par
}

\let\ResumeDemonstrations\Demonstrations\renewcommand{\Demonstrations}{}\renewcommand{\OtherExperience}{}
